% Chapter 8, Section 4

\section{Second-Order Methods}
\label{sec:second-order}

\subsection{Intuition: Curvature-Aware Steps}

First-order methods follow the slope; second-order methods also consider the \textit{curvature} of the landscape to choose better-scaled steps. If the valley is sharply curved in one direction and flat in another, curvature-aware steps shorten strides along the sharp direction and lengthen them along the flat one.\index{second-order methods}

Historical note: Classical optimization popularized Newton and quasi-Newton methods; in deep learning, memory and compute constraints motivated approximations like L-BFGS and natural gradient \cite{LiuNocedal1989,Amari1998,GoodfellowEtAl2016,Bishop2006}.

\subsection{Newton's Method}

Uses second-order Taylor expansion:
\begin{equation}
\vect{\theta}_{t+1} = \vect{\theta}_t - \mat{H}^{-1} \nabla_{\vect{\theta}} L(\vect{\theta}_t)
\end{equation}

where $\mat{H}$ is the Hessian matrix.

Challenges:\index{Newton's method}
\begin{itemize}
    \item Computing Hessian is $O(n^2)$ in parameters
    \item Inverting Hessian is $O(n^3)$
    \item Infeasible for large neural networks
\end{itemize}

\subsection{Quasi-Newton Methods}\index{L-BFGS}

Approximate the Hessian inverse:

\textbf{L-BFGS:} maintains low-rank approximation of Hessian inverse
\begin{itemize}
    \item More efficient than full Newton's method
    \item Still expensive for very large models
    \item Used for smaller networks or specific applications
\end{itemize}

\subsection{Natural Gradient}\index{natural gradient}

Uses Fisher information matrix instead of Hessian:
\begin{equation}
\vect{\theta}_{t+1} = \vect{\theta}_t - \alpha \mat{F}^{-1} \nabla_{\vect{\theta}} L(\vect{\theta}_t)
\end{equation}

Provides parameter updates invariant to reparameterization.

